% Section 2.8, from lectures/L02/appendix/c_testbenches.md: the course's guide to running
% testbenches, one subsection per part, C.1 to C.7. Every later chapter points back here.

\section{Verifying Your VHDL with a Testbench}
\label{c2:app:c}
This section is the book's guide to \textbf{running} testbenches, and every later chapter's
exercises point back here. You do not write testbenches in this book, but you run one after every
VHDL exercise, because with no FPGA board of your own that is what verifies your VHDL. Every worked example and every exercise ships a
ready-made \file{<module>_tb.vhd}.

\subsection{What you need}
\label{c2:c1}
\textbf{GHDL}, the open-source VHDL analyzer and simulator. Check with \code{ghdl --version}; install
on WSL/Ubuntu with \code{sudo apt -y install ghdl}.

Every command below uses \code{--std=93}, the VHDL standard this book targets.

\subsection{What a testbench is}
\label{c2:c2}
A testbench is itself a VHDL module, but a special one:
\begin{itemize}
  \item Its entity has \textbf{no ports}: nothing connects to the outside, because it is the top of
    the simulation.
  \item Its architecture \textbf{instantiates your design}, the ``device under test'' labelled
    \code{dut}, and drives its inputs.
  \item It \textbf{checks} each output with an \code{assert}. If an output is wrong, the assertion
    fires with a message and the simulation stops.
\end{itemize}
So instead of you reading a waveform to decide whether the output is right, the testbench decides,
and only speaks up when something is wrong.

\subsection{Running a testbench: analyze \texorpdfstring{$\to$}{->} elaborate
  \texorpdfstring{$\to$}{->} run}
\label{c2:c3}
Three steps. Pass \code{--std=93} to all three, and give the \textbf{testbench entity name}, not a
filename, to the last two:

\begin{shell}
ghdl -a --std=93 <module>.vhd <module>_tb.vhd  # analyze both files (module first)
ghdl -e --std=93 <module>_tb                   # elaborate the testbench
ghdl -r --std=93 <module>_tb --assert-level=error --stop-time=10ms   # run it
\end{shell}

\begin{itemize}
  \item \code{--assert-level=error} makes a failed check stop the run and return a
    \textbf{non-zero exit code}, a real pass/fail you can check with \code{echo $?}. The testbenches
    here raise failures at \code{severity failure}, which stops the run on its own, so you get a
    real pass/fail even without the flag. Keep it on anyway: it is the habit that also works with
    testbenches reporting at a lower severity.
  \item \code{--stop-time=10ms} caps how long the simulation may run. A testbench often waits for
    your module to do something, with a line like \code{wait until tx = '0';}. If your module never
    does, that wait never finishes and without a stop time the simulation hangs: no output, no
    error, just a terminal sitting there. 10\,ms of \emph{simulated} time is far more than any
    testbench here needs, and takes well under a second of real time.
  \item A passing run prints its ``all checks passed'' note and exits 0. Every testbench in this
    book ends with that note, so if you do not see it, the run did not finish: it either failed a
    check or hit the stop time.
\end{itemize}

\subsection{Checking an exercise}
\label{c2:c4}
Each VHDL exercise ships its testbench in its own directory. For example, this chapter's
Exercise~\ref{c2:ex:3} has \file{lectures/L02/exercises/xyz_logic/xyz_logic_tb.vhd}. Write your
module in \textbf{that same directory}, named exactly as the exercise asks, then run the three
commands from there:

\begin{shell}
cd lectures/L02/exercises/xyz_logic
ghdl -a --std=93 xyz_logic.vhd xyz_logic_tb.vhd
ghdl -e --std=93 xyz_logic_tb
ghdl -r --std=93 xyz_logic_tb --assert-level=error --stop-time=10ms
\end{shell}

\noindent Each testbench header lists these exact commands.

The testbench connects to your module \textbf{by position}, not by name, so your entity must declare
its ports in the \textbf{same order} the exercise lists them. The names are yours to choose, but use
the ones specified anyway, so your module and the testbench's error messages talk about the same
signals.

A port declared in the wrong order fails in one of two ways, neither of which says ``wrong order'':
\begin{itemize}
  \item if the misplaced ports have \textbf{different} types or widths, analysis stops with
    \code{can't associate "sel" with port "sel"}, pointing at the testbench's \code{port map} line.
  \item if they have the \textbf{same} type, several \code{std_logic} ports say, it analyzes and
    elaborates quite happily, then fails at run time as an ordinary assertion failure, which reads
    exactly like a logic bug in your design.
\end{itemize}
So when a testbench fails on the very first case it checks, re-read your port order before hunting
through your architecture. The same applies to \textbf{generics}: a module that carries them must
declare exactly the ones the exercise specifies, in order.

\subsection{When the design needs more than one file}
\label{c2:c5}
Some exercises build a module out of another one, so the \code{ghdl -a} line names several files
instead of two. The elaborate and run steps are unchanged: they still take the testbench
\textbf{entity name} only.

\needspace{5\baselineskip}
\begin{shell}
ghdl -a --std=93 <subcomponent>.vhd <module>.vhd <module>_tb.vhd
ghdl -e --std=93 <module>_tb
ghdl -r --std=93 <module>_tb --assert-level=error --stop-time=10ms
\end{shell}

\begin{itemize}
  \item \textbf{Order matters on the analyze line}: name a module before anything that instantiates
    it, so GHDL has already seen the entity when it reaches the \code{entity work.<name>} that
    refers to it. Subcomponents first, your top level next, the testbench last.
  \item No exercise directory ships a subcomponent. Every module a design instantiates is one you
    wrote in an earlier exercise and copy in yourself; the exercise says which, and where from.
  \item The exercises that need this are this chapter's \code{hex_display}
    (Exercise~\ref{c2:ex:8}, which instantiates your \code{display}), Chapter~\ref{c4:ch}'s
    \code{led_toggle_sync2}, Chapter~\ref{c7:ch}'s \code{blinker} (\code{reset_sync} and
    \code{timer}) and \code{walking_led} (those two plus \code{button_sync}), and all of
    Chapter~\ref{c8:ch}'s, where the receiver names five files and the transmitter four.
  \item Every testbench header lists the exact command for its own exercise.
\end{itemize}

\subsection{Reading the result}
\label{c2:c6}
A failed check looks like this:

\begin{console}
xyz_logic_tb.vhd:37:13:@10ns:(assertion failure): xyz_logic: wrong X with a='0' b='0' c='0' d='0', expected '0' but got '1'!
ghdl:error: assertion failed
\end{console}

\noindent Read it in order: \code{file:line} is where the assert is, \code{@time} is when it fired,
then the message. It is the same shape as a GHDL compile error, so it maps onto what you already
know. No failure means every case matched.

\subsection{A complete example}
\label{c2:c7}
A combinational testbench in the shape every testbench in this book uses, here for
Chapter~\ref{c1:ch}'s \code{or_gate}:

\begin{vhdlcode}
library ieee;
use ieee.std_logic_1164.all;
use ieee.numeric_std.all;

entity or_gate_tb is
end entity;

architecture behaviour of or_gate_tb is
signal a, b, x: std_logic;
begin
    dut: entity work.or_gate
        port map(a, b, x);   -- positional: same order as or_gate's port clause.

    SIM_PROCESS: process is
    begin
        for i in 0 to 3 loop
            (a, b) <= std_logic_vector(to_unsigned(i, 2)); -- apply inputs
            wait for 10 ns;                                -- let them settle
            assert x = (a or b)                            -- check AFTER the wait
                report "or_gate: wrong output with a=" & std_logic'image(a)
                     & " b=" & std_logic'image(b)
                     & ", expected " & std_logic'image(a or b)
                     & " but got " & std_logic'image(x) & "!"
                severity failure;
        end loop;
        report "or_gate: all checks passed!" severity note;
        wait;                                              -- end the simulation
    end process;
end architecture;
\end{vhdlcode}

\noindent Two rules keep a combinational testbench correct:
\begin{itemize}
  \item \textbf{Apply $\to$ wait $\to$ check, in that order.} A signal assignment only
    \emph{schedules} a value, so the inputs and output do not update until the \code{wait}.
    Checking before it tests the previous iteration's values.
  \item \textbf{End with \code{wait;}} so the stimulus process suspends and the simulation stops on
    its own.
\end{itemize}
Clocked designs need a little more, a clock-generator process and a way to stop it, which is exactly
the machinery the provided testbenches already contain. For the exercises you only ever write the
module.
